\documentclass[a4paper,11pt]{article}
\usepackage[T2A]{fontenc}
\usepackage[utf8]{inputenc}
\usepackage[russian]{babel}
\usepackage{amsmath,amssymb,mathtools}
\usepackage{geometry}
\geometry{left=21mm,right=21mm,top=20mm,bottom=20mm,headheight=22pt}
\usepackage{microtype}
\usepackage{booktabs,tabularx,array}
\usepackage{enumitem}
\usepackage{tikz}
\usepackage{fancyhdr}
\usepackage{hyperref}
\usepackage{setspace}
\usepackage{xcolor}
\definecolor{accent}{HTML}{087D83}
\definecolor{darkaccent}{HTML}{075A61}
\definecolor{ink}{HTML}{303738}
\definecolor{soft}{HTML}{6C777A}
\definecolor{error}{HTML}{A53332}
\hypersetup{colorlinks=true,linkcolor=darkaccent,urlcolor=accent}
\setstretch{1.06}
\setlength{\parindent}{0pt}
\setlength{\parskip}{4pt}
\setlist[itemize]{topsep=3pt,itemsep=3pt,leftmargin=1.4em}
\pagestyle{fancy}
\fancyhf{}
\fancyhead[L]{\small\color{soft}Ревью домашней работы}
\fancyhead[R]{\small\color{soft}Анастасия Павлова}
\renewcommand{\headrulewidth}{0.45pt}
\renewcommand{\headrule}{\hbox to\headwidth{\color{accent}\leaders\hrule height \headrulewidth\hfill}}
\fancyfoot[C]{\small\color{soft}\thepage}
\newcommand{\topic}[1]{\vspace{5pt}{\large\color{darkaccent}#1}\par\vspace{2pt}}
\newcommand{\err}{\textcolor{error}{\textbf{Ошибка:}}}
\newcommand{\idea}{\textcolor{darkaccent}{\textbf{Ключевая идея:}}}
\newcommand{\result}{\textcolor{darkaccent}{\textbf{Итог:}}}
\begin{document}
\thispagestyle{empty}
\begin{center}
\vspace*{19mm}
{\large\color{accent} МАТЕМАТИКА}\par
\vspace{30mm}
{\Huge Ревью домашней работы}\par
\vspace{15mm}
{\Large\color{accent} Кредиты и равные ежегодные платежи}\par
\vspace{14mm}
{\normalsize Проценты. Коэффициент роста. Погашение долга.}\par
\vfill
{\large Анастасия Павлова}\par
\vspace{10mm}
{5 ноября 2026 г.}\par
\vspace{19mm}
{\small Лёвин Артём Александрович}\par
{\small эксклюзивно для Анастасии Павловой}\par
\vspace{15mm}
\begin{tikzpicture}[x=1cm,y=1cm]
\draw[accent,line width=.65pt] (0,0) .. controls (1.0,.45) and (2.5,-.35) .. (3.6,0)
 .. controls (4.7,.4) and (6,-.45) .. (7.4,0);
\end{tikzpicture}
\end{center}
\clearpage
\section*{\color{darkaccent}Сводная таблица}
В таблицу включены только задания, требующие исправления. Номер задания ведёт к подробному разбору.
\par\vspace{9mm}
\noindent{\footnotesize
\setlength{\tabcolsep}{3pt}
\renewcommand{\arraystretch}{1.5}
\begin{tabularx}{\linewidth}{@{}p{11mm}p{24mm}X p{27mm}p{28mm}@{}}
\toprule
\textbf{№} & \textbf{Тема} & \textbf{Суть ошибки} & \textbf{Ваш ответ} & \textbf{Правильный ответ}\\
\midrule
\hyperref[task:7]{\textbf{7}} & Трёхлетний кредит & Применена схема двух платежей; неверно учтено начисление процентов. & $\approx118\,054$ руб. & $86\,400$ руб.\\
\addlinespace[6pt]
\hyperref[task:8]{\textbf{8}} & Двухлетний кредит & Неверный множитель при первом платеже; вместо двух платежей в сумме учтены три. & Переплата $104\,053$ руб. & Переплата $34\,000$ руб.\\
\bottomrule
\end{tabularx}}
\vspace{15mm}
\topic{Общий принцип}
Если долг увеличивается за год в $k$ раз, а в конце каждого года вносится одинаковый платёж $x$, то для кредита на $n$ лет:
\[
 Sk^n-x\bigl(k^{n-1}+k^{n-2}+\cdots+k+1\bigr)=0.
\]
Здесь $S$ --- первоначальная сумма кредита. Каждый платёж учитывается ровно один раз; чем раньше он внесён, тем больше раз уменьшенная на него задолженность участвует в последующих начислениях процентов.
\clearpage
\thispagestyle{empty}
\vspace*{\fill}
\begin{center}
{\Huge\color{darkaccent}Разбор ошибок}
\end{center}
\vspace*{\fill}
\clearpage
\phantomsection\label{task:7}
\section*{\color{darkaccent}Задание № 7. Кредит на три года}
\topic{Текст задания}
Кредит $182\,000$ руб. взят на три года под $20\%$ годовых. Найдите ежегодный равный платёж.

\topic{Ваше решение}
В записи составлено уравнение с двумя вычитаемыми платежами:
\[
 182\cdot1{,}44-1{,}22x-x=0,
\]
после чего получено приближённое значение $x\approx118{,}054$ тыс. руб., то есть около $118\,054$ руб.

\topic{Ваш ответ}
$118\,054$ руб. (приближённо).

\topic{Ошибка}
\err\ Применена двухлетняя схема к трёхлетнему кредиту. Кроме того, коэффициент при первом платеже записан как $1{,}22$; для годовой ставки $20\%$ соответствующий коэффициент за один год равен $1{,}2$.

\topic{Где именно возникает ошибка}
Последний верный шаг --- определение годового коэффициента:
\[
 k=1+\frac{20}{100}=1{,}2.
\]
Первый неверный шаг --- запись $182\cdot1{,}44-1{,}22x-x=0$.
Коэффициент $1{,}44=1{,}2^2$ описывает рост за два года; в условии кредит погашается \emph{за три года}. Поэтому необходимо учитывать рост первоначального долга за три года и \emph{три} платежа.

\topic{Почему это неверно}
В конце первого года вносится первый платёж. Он уменьшает долг, который затем ещё два года подвергался бы начислению процентов. Второй платёж вносится в конце второго года и влияет на начисление процентов в течение одного оставшегося года. Третий платёж вносится в конце третьего года и сразу завершает расчёт. Таким образом, при приведении всех сумм к концу третьего года коэффициенты трёх платежей равны $1{,}2^2$, $1{,}2$ и $1$.

\topic{Ключевая идея}
\idea\ Для кредита на три года сравнивайте задолженность и платежи \emph{на одну и ту же дату} --- конец третьего года.
\[
 182\cdot1{,}2^3-x\cdot1{,}2^2-x\cdot1{,}2-x=0.
\]

\topic{Исправление от последнего верного шага}
Используем найденное $k=1{,}2$, записываем три ежегодных платежа и собираем слагаемые с $x$:
\[
 182\cdot1{,}2^3-x(1{,}2^2+1{,}2+1)=0.
\]
\clearpage
\section*{\color{darkaccent}Продолжение разбора}
\topic{Полное правильное решение}
Для удобства все суммы будем выражать в тысячах рублей. Пусть $x$ --- ежегодный платёж в тысячах рублей.

В конце первого года после начисления процентов и платежа остаётся:
\[
 D_1=182\cdot1{,}2-x=218{,}4-x.
\]
К концу второго года, после следующего начисления процентов и второго платежа:
\begin{align*}
D_2&=(218{,}4-x)\cdot1{,}2-x\\
   &=262{,}08-1{,}2x-x\\
   &=262{,}08-2{,}2x.
\end{align*}
После третьего года задолженность должна стать нулевой:
\begin{align*}
0&=(262{,}08-2{,}2x)\cdot1{,}2-x,\\
0&=314{,}496-2{,}64x-x,\\
3{,}64x&=314{,}496,\\
x&=\frac{314{,}496}{3{,}64}=86{,}4.
\end{align*}
Следовательно, ежегодный платёж равен $86{,}4$ тыс. руб.

\topic{Проверка}
Проверим остаток после каждого из трёх платежей (тыс. руб.):
\begin{align*}
\text{Год 1:}\quad&182\cdot1{,}2-86{,}4=132;\\
\text{Год 2:}\quad&132\cdot1{,}2-86{,}4=72;\\
\text{Год 3:}\quad&72\cdot1{,}2-86{,}4=0.
\end{align*}
После третьего платежа долг погашен полностью. Все три платежа одинаковы.

\topic{Итог}
\result\ Ежегодный равный платёж --- \textbf{$86\,400$ руб.}

\topic{Задача на закрепление}
Кредит $91\,000$ руб. взят на три года под $10\%$ годовых. Кредит погашается тремя равными платежами в конце каждого года. Найдите размер ежегодного платежа.
\clearpage
\phantomsection\label{task:8}
\section*{\color{darkaccent}Задание № 8. Переплата по кредиту}
\topic{Текст задания}
Кредит $110\,000$ руб. под $20\%$ годовых погашают двумя равными ежегодными платежами. Найдите общую сумму платежей и переплату в рублях.

\topic{Ваше решение}
После верного вычисления $110\cdot1{,}44=158{,}4$ (тыс. руб.) записано:
\[
158{,}4-1{,}22x-x=0,\qquad
x\approx71{,}351\ \text{(тыс. руб.)}.
\]
Далее ежегодный платёж умножен на три:
\[
71\,351\cdot3=214\,053\ \text{руб.},
\]
и из полученной суммы вычтено $110\,000$ руб.

\topic{Ваш ответ}
Общая сумма платежей --- $214\,053$ руб.; переплата --- $104\,053$ руб.

\topic{Ошибка}
\err\ В уравнении коэффициент при первом платеже записан $1{,}22$ вместо $1{,}2$. При подсчёте итоговой суммы использованы три платежа, хотя по условию их два.

\topic{Где именно возникает ошибка}
Последний правильный шаг --- начисление процентов на исходную сумму за два года:
\[
110\cdot1{,}2^2=110\cdot1{,}44=158{,}4.
\]
Первый неправильный шаг --- уравнение $158{,}4-1{,}22x-x=0$: первый платёж после внесения влияет на задолженность в течение \emph{одного} следующего года, поэтому умножается на $1{,}2$.

\topic{Почему это неверно}
В конце первого года выплачивается $x$. Если перенести влияние этого платежа к концу второго года, нужно умножить его на один годовой коэффициент $1{,}2$. Запись $1{,}22$ не следует из процентной ставки. В конце второго года вносится ещё один платёж $x$. Кредит рассчитан на два платежа; следовательно, общая сумма выплат равна $2x$.

\topic{Ключевая идея}
\idea\ Сначала правильно рассчитайте одинаковый платёж, затем умножьте его на число платежей, прямо указанное в условии.
\[
110\cdot1{,}2^2-1{,}2x-x=0.
\]
\topic{Исправление от последнего верного шага}
Сохраняем верно найденные $158{,}4$ тыс. руб. и исправляем коэффициент:
\[
158{,}4-1{,}2x-x=0.
\]

\clearpage
\section*{\color{darkaccent}Продолжение разбора}
\topic{Полное правильное решение}
Пусть $x$ --- ежегодный платёж в тысячах рублей. Ставка $20\%$ означает, что за год задолженность увеличивается в $1{,}2$ раза.

Первый год. До платежа задолженность:
\[
110\cdot1{,}2=132.
\]
После внесения первого платежа остаётся $132-x$.

Второй год. Снова начисляются проценты, затем вносится второй платёж. По условию долг погашен:
\begin{align*}
(132-x)\cdot1{,}2-x&=0,\\
158{,}4-1{,}2x-x&=0,\\
2{,}2x&=158{,}4,\\
x&=\frac{158{,}4}{2{,}2}=72.
\end{align*}
Каждый платёж равен $72$ тыс. руб., то есть $72\,000$ руб.

Всего платежей два:
\[
72\,000\cdot2=144\,000\ \text{руб.}
\]
Переплата --- разность между всей выплаченной суммой и первоначальной суммой кредита:
\[
144\,000-110\,000=34\,000\ \text{руб.}
\]

\topic{Проверка}
После первого года:
\[
110\cdot1{,}2-72=60\ \text{(тыс. руб.)}.
\]
После второго года:
\[
60\cdot1{,}2-72=0.
\]
Задолженность стала равной нулю ровно после двух одинаковых платежей.

\topic{Итог}
\result\ Общая сумма платежей --- \textbf{$144\,000$ руб.}; переплата --- \textbf{$34\,000$ руб.}

\topic{Задача на закрепление}
Кредит $80\,000$ руб. под $25\%$ годовых погашают двумя равными платежами в конце каждого года. Найдите общую сумму платежей и переплату в рублях.
\end{document}